
\documentclass{article}
\usepackage{colortbl}
\usepackage{makecell}
\usepackage{multirow}
\usepackage{supertabular}

\begin{document}

\newcounter{utterance}

\centering \large Interaction Transcript for game `hot\_air\_balloon', experiment `air\_balloon\_survival\_lt\_negotiation\_easy', episode 0 with Qwen3.6{-}35B{-}A3B{-}FP8{-}without{-}reasoning.
\vspace{24pt}

{ \footnotesize  \setcounter{utterance}{1}
\setlength{\tabcolsep}{0pt}
\begin{supertabular}{c@{$\;$}|p{.15\linewidth}@{}p{.15\linewidth}p{.15\linewidth}p{.15\linewidth}p{.15\linewidth}p{.15\linewidth}}
    \# & $\;$A & \multicolumn{4}{c}{Game Master} & $\;\:$B\\
    \hline

    \theutterance \stepcounter{utterance}  
    & & \multicolumn{4}{p{0.6\linewidth}}{
        \cellcolor[rgb]{0.9,0.9,0.9}{
            \makecell[{{p{\linewidth}}}]{
                \texttt{\tiny{[P1$\langle$GM]}}
                \texttt{Dalyvaujate bendradarbiavimo derybų žaidime.} \\
\\ 
\texttt{Kartu su kitu dalyviu turite susitarti dėl vieno daiktų rinkinio, kuris bus išlaikytas. Kiekvienas iš jūsų turi savo nuomonę apie tai, kiek kiekvienas daiktas jums svarbus (svarba). Jūs nežinote, kaip kitas dalyvis vertina daiktus. Be to, jums nurodoma, kiek pastangų reikalauja kiekvienas daiktas.} \\
\texttt{Dėl rinkinio galite susitarti tik tuo atveju, jei pasirinktų daiktų bendros pastangos neviršija bendros ribos:} \\
\\ 
\texttt{LIMIT = 3358} \\
\\ 
\texttt{Čia pateikiamos atskirų daiktų pastangų vertės:} \\
\\ 
\texttt{Daikto pastangos = \{"C08": 171, "C55": 933, "A43": 552, "A67": 219, "B42": 275, "A62": 778, "B19": 341, "C43": 615, "C73": 519, "C10": 862, "B36": 262, "A18": 377, "C61": 347, "C67": 349, "C42": 117\}} \\
\\ 
\texttt{Čia pateikiama jūsų asmeninė nuomonė apie kiekvieno daikto svarbą:} \\
\\ 
\texttt{Daikto svarbos reikšmės = \{"C08": 138, "C55": 583, "A43": 868, "A67": 822, "B42": 783, "A62": 65, "B19": 262, "C43": 121, "C73": 508, "C10": 780, "B36": 461, "A18": 484, "C61": 668, "C67": 389, "C42": 808\}} \\
\\ 
\texttt{Tikslas:} \\
\\ 
\texttt{Jūsų tikslas {-} susitarti dėl bendro daiktų rinkinio, kuris būtų jums kuo naudingesnis (t. y. maksimaliai padidintų bendrą svarbą JUMS), neviršijant LIMIT. Neprivalote kiekvienoje žinutėje pateikti PASIŪLYMAS {-} galite tiesiog derėtis. Visos taktikos yra leidžiamos!} \\
\\ 
\texttt{Sąveikos protokolas:} \\
\\ 
\texttt{Žinutėje galite naudoti tik toliau nurodytus struktūrinius formatus:} \\
\\ 
\texttt{PASIŪLYMAS: \{'A', 'B', 'C', …\}} \\
\texttt{Siūlykite išlaikyti būtent šiuos daiktus.} \\
\texttt{ATSISAKAU: \{'A', 'B', 'C', …\}} \\
\texttt{Aiškiai atmeskite oponento pasiūlymą.} \\
\texttt{ARGUMENTAS: \{'...'\}} \\
\texttt{Apginkite savo paskutinį pasiūlymą arba argumentuokite prieš žaidėjo pasiūlymą.} \\
\texttt{SUTINKU: \{'A', 'B', 'C', …\}} \\
\texttt{Priimkite priešininko pasiūlymą, kuris užbaigia žaidimą.} \\
\\ 
\\ 
\texttt{Taisyklės:} \\
\\ 
\texttt{Galite SUTINKU tik su pasiūlymu, kurį kita šalis užregistravo per PASIŪLYMAS.} \\
\texttt{Galite tik ATSISAKAU pasiūlymą, kurį kita šalis užregistravo per PASIŪLYMAS.} \\
\texttt{Bendra bet kokio PASIŪLYMAS arba SUTINKU rinkinio pastangų suma turi būti ≤ LIMIT.} \\
\texttt{NESKELBKITE savo paslėptų svarbos balų.} \\
\texttt{Po struktūrinio formato žymos turi sekti dvitaškis ir tarpas. Argumentas turi būti python rinkinys, kuriame yra 0 ar daugiau eilučių.} \\
\texttt{Taigi jis turi būti tokios formos: "TAG: \{...\}"} \\
\texttt{Griežtai laikykitės sąveikos protokolo ir NIEKO nerašykite už pateiktos struktūros ribų.} \\
\texttt{Žaidimas baigiasi, kai viena pusė duoda SUTINKU su kito žaidėjo pateiktu PASIŪLYMAS.} \\
\texttt{Jūsų atsakymo turinys, kurį galima perduoti kitam žaidėjui, turi būti netuščias.} \\
\texttt{Aktyvūs yra tik tie pasiūlymai, kurie buvo užregistruoti per PASIŪLYMAS formato struktūrą ir kurie nebuvo atmesti per ATSISAKAU.} \\
\\ 
            }
        }
    }
    & \\ \\

    \theutterance \stepcounter{utterance}  
    & \multicolumn{4}{p{0.6\linewidth}}{
        \cellcolor[rgb]{0.9,0.9,0.9}{
            \makecell[{{p{\linewidth}}}]{
                \texttt{\tiny{[P1$\rangle$GM]}}
                \texttt{PASIŪLYMAS: \{'C08', 'A67', 'B42', 'B36', 'C42', 'C61', 'C67', 'B19', 'C10'\}} \\
            }
        }
    }
    & & \\ \\

    \theutterance \stepcounter{utterance}  
    & & \multicolumn{4}{p{0.6\linewidth}}{
        \cellcolor[rgb]{0.9,0.9,0.9}{
            \makecell[{{p{\linewidth}}}]{
                \texttt{\tiny{[GM$\rangle$P2]}}
                \texttt{Dalyvaujate bendrų derybų žaidime.} \\
\\ 
\texttt{Kartu su kitu dalyviu turite susitarti dėl vieno daiktų rinkinio, kuris bus išsaugotas. Kiekvienas iš jūsų turi savo nuomonę, kiek kiekvienas daiktas jums svarbus (svarba). Jūs nežinote, kaip daiktus vertina kitas dalyvis. Be to, jums nurodoma, kiek pastangų reikalauja kiekvienas daiktas.} \\
\texttt{Dėl rinkinio galite susitarti tik tuo atveju, jei bendra pasirinktų daiktų pastangų suma neviršija bendros ribos:} \\
\\ 
\texttt{LIMIT = 3358} \\
\\ 
\texttt{Čia pateikiamos atskirų daiktų pastangų vertės:} \\
\\ 
\texttt{Daikto pastangos = \{"C08": 171, "C55": 933, "A43": 552, "A67": 219, "B42": 275, "A62": 778, "B19": 341, "C43": 615, "C73": 519, "C10": 862, "B36": 262, "A18": 377, "C61": 347, "C67": 349, "C42": 117\}} \\
\\ 
\texttt{Čia pateikiama jūsų asmeninė nuomonė apie kiekvieno daikto svarbą:} \\
\\ 
\texttt{Daikto svarbos reikšmės = \{"C08": 138, "C55": 583, "A43": 868, "A67": 822, "B42": 783, "A62": 65, "B19": 262, "C43": 121, "C73": 508, "C10": 780, "B36": 461, "A18": 484, "C61": 668, "C67": 389, "C42": 808\}} \\
\\ 
\texttt{Tikslas:} \\
\\ 
\texttt{Jūsų tikslas {-} susitarti dėl bendro daiktų rinkinio, kuris būtų jums kuo naudingesnis (t. y. maksimaliai padidintų bendrą svarbą JUMS), neviršijant LIMITO. Neprivalote kiekvienoje žinutėje pateikti PASIŪLYMO {-} galite tiesiog derėtis taip pat. Visos taktikos yra leidžiamos!} \\
\\ 
\texttt{Sąveikos protokolas:} \\
\\ 
\texttt{Žinutėje galite naudoti tik toliau nurodytus struktūrinius formatus:} \\
\\ 
\texttt{PASIŪLYMAS: \{'A', 'B', 'C', …\}} \\
\texttt{Siūlykite išlaikyti būtent šiuos daiktus.} \\
\texttt{ATSISAKAU: \{'A', 'B', 'C', …\}} \\
\texttt{Aiškiai atmesti oponento pasiūlymą.} \\
\texttt{ARGUMENTAS: \{'...'\}} \\
\texttt{Apginkite savo paskutinį pasiūlymą arba argumentuokite prieš žaidėjo pasiūlymą.} \\
\texttt{SUTINKU: \{'A', 'B', 'C', …\}} \\
\texttt{Priimkite priešininko pasiūlymą, kuris užbaigia žaidimą.} \\
\\ 
\\ 
\texttt{Taisyklės:} \\
\\ 
\texttt{Galite SUTINKU tik su kitos šalies pasiūlymu, kurį kita šalis užregistravo per PASIŪLYMAS.} \\
\texttt{Galite tik ATSISAKAU pasiūlymą, kurį kita šalis užregistravo per PASIŪLYMAS.} \\
\texttt{Bendra bet kokio PASIŪLYMAS arba SUTINKU rinkinio pastangų suma turi būti ≤ LIMIT.} \\
\texttt{NESKELBKITE savo paslėptų svarbos balų.} \\
\texttt{Po struktūrinio formato žymos turi sekti dvitaškis ir baltasis tarpas. Argumentas turi būti python rinkinys, kuriame yra 0 ar daugiau eilučių.} \\
\texttt{Taigi ji turi būti tokios formos: "TAG: \{...\}"} \\
\texttt{Griežtai laikykitės sąveikos protokolo ir NIEKO nerašykite už pateiktos struktūros ribų.} \\
\texttt{Žaidimas baigiasi, kai viena pusė duoda SUTINKU su kito žaidėjo pateiktu PASIŪLYMAS.} \\
\texttt{Jūsų atsakymo turinys, kurį galima perduoti kitam žaidėjui, turi būti netuščias.} \\
\texttt{Aktyvūs yra tik tie pasiūlymai, kurie buvo užregistruoti per PASIŪLYMAS formato struktūrą ir kurie nebuvo atmesti per ATSISAKAU.} \\
\\ 
\\ 
\\ 
\texttt{Dabar gausite pirmąjį kito žaidėjo pranešimą.} \\
\\ 
\texttt{PASIŪLYMAS: \{'C08', 'A67', 'B42', 'B36', 'C42', 'C61', 'C67', 'B19', 'C10'\}} \\
            }
        }
    }
    & \\ \\

    \theutterance \stepcounter{utterance}  
    & & & \multicolumn{4}{p{0.6\linewidth}}{
        \cellcolor[rgb]{0.9,0.9,0.9}{
            \makecell[{{p{\linewidth}}}]{
                \texttt{\tiny{[GM$\langle$P2]}}
                \texttt{PASIŪLYMAS: \{'A43', 'A67', 'B42', 'B19', 'C61', 'C67', 'C73', 'C43', 'A18', 'C08', 'B36'\}} \\
            }
        }
    }
     \\ \\

    \theutterance \stepcounter{utterance}  
    & & \multicolumn{4}{p{0.6\linewidth}}{
        \cellcolor[rgb]{0.9,0.9,0.9}{
            \makecell[{{p{\linewidth}}}]{
                \texttt{\tiny{[P1$\langle$GM]}}
                \texttt{PASIŪLYMAS: \{'A43', 'A67', 'B42', 'B19', 'C61', 'C67', 'C73', 'C43', 'A18', 'C08', 'B36'\}} \\
            }
        }
    }
    & \\ \\

    \theutterance \stepcounter{utterance}  
    & \multicolumn{4}{p{0.6\linewidth}}{
        \cellcolor[rgb]{0.9,0.9,0.9}{
            \makecell[{{p{\linewidth}}}]{
                \texttt{\tiny{[P1$\rangle$GM]}}
                \texttt{ATSISAKAU: \{'A43', 'A67', 'B42', 'B19', 'C61', 'C67', 'C73', 'C43', 'A18', 'C08', 'B36'\}} \\
\\ 
\texttt{ARGUMENTAS: \{'Jūsų pasiūlyme yra daiktas C43 (pastangos: 615), kuris man neturi jokios vertės (svarba: 121). Vietoj jo aš siūlau C10, kuris man yra daug vertingesnis (svarba: 780), o pastangos panašios (862 vs 615). Be to, išsaugome vietą biudžete. Mano siūlomas rinkinis suteikia man didesnę bendrą naudą. Bandykime susitarti dėl balanso, kuris gerbtų mūsų abiejų interesus.'\}} \\
            }
        }
    }
    & & \\ \\

    \theutterance \stepcounter{utterance}  
    & & \multicolumn{4}{p{0.6\linewidth}}{
        \cellcolor[rgb]{0.9,0.9,0.9}{
            \makecell[{{p{\linewidth}}}]{
                \texttt{\tiny{[GM$\rangle$P2]}}
                \texttt{ATSISAKAU: \{'A43', 'A67', 'B42', 'B19', 'C61', 'C67', 'C73', 'C43', 'A18', 'C08', 'B36'\}} \\
\\ 
\texttt{ARGUMENTAS: \{'Jūsų pasiūlyme yra daiktas C43 (pastangos: 615), kuris man neturi jokios vertės (svarba: 121). Vietoj jo aš siūlau C10, kuris man yra daug vertingesnis (svarba: 780), o pastangos panašios (862 vs 615). Be to, išsaugome vietą biudžete. Mano siūlomas rinkinis suteikia man didesnę bendrą naudą. Bandykime susitarti dėl balanso, kuris gerbtų mūsų abiejų interesus.'\}} \\
            }
        }
    }
    & \\ \\

    \theutterance \stepcounter{utterance}  
    & & & \multicolumn{4}{p{0.6\linewidth}}{
        \cellcolor[rgb]{0.9,0.9,0.9}{
            \makecell[{{p{\linewidth}}}]{
                \texttt{\tiny{[GM$\langle$P2]}}
                \texttt{ATSISAKAU: \{'C08', 'A67', 'B42', 'B36', 'C42', 'C61', 'C67', 'B19', 'C10'\}} \\
\\ 
\texttt{ARGUMENTAS: \{'Mano pirmajame pasiūlyme svarbiausi daiktai buvo C10 (svarba 780, pastangos 862), A43 (svarba 868, pastangos 552) ir C42 (svarba 808, pastangos 117). Jūsų atsakyme šie kritiniai elementai atmetami. Jūsų siūlomas C43 (pastangos 615, svarba 121) yra neprotingas pasirinkimas abiem pusėms, nes jo nauda minimali, o kaštai dideli. Siūlau naują balansuotą variantą, įtraukiantį C42, A43 ir C10, kurie yra didžiausios naudos šaltiniai.'\}} \\
            }
        }
    }
     \\ \\

    \theutterance \stepcounter{utterance}  
    & & \multicolumn{4}{p{0.6\linewidth}}{
        \cellcolor[rgb]{0.9,0.9,0.9}{
            \makecell[{{p{\linewidth}}}]{
                \texttt{\tiny{[P1$\langle$GM]}}
                \texttt{ATSISAKAU: \{'C08', 'A67', 'B42', 'B36', 'C42', 'C61', 'C67', 'B19', 'C10'\}} \\
\\ 
\texttt{ARGUMENTAS: \{'Mano pirmajame pasiūlyme svarbiausi daiktai buvo C10 (svarba 780, pastangos 862), A43 (svarba 868, pastangos 552) ir C42 (svarba 808, pastangos 117). Jūsų atsakyme šie kritiniai elementai atmetami. Jūsų siūlomas C43 (pastangos 615, svarba 121) yra neprotingas pasirinkimas abiem pusėms, nes jo nauda minimali, o kaštai dideli. Siūlau naują balansuotą variantą, įtraukiantį C42, A43 ir C10, kurie yra didžiausios naudos šaltiniai.'\}} \\
            }
        }
    }
    & \\ \\

    \theutterance \stepcounter{utterance}  
    & \multicolumn{4}{p{0.6\linewidth}}{
        \cellcolor[rgb]{0.9,0.9,0.9}{
            \makecell[{{p{\linewidth}}}]{
                \texttt{\tiny{[P1$\rangle$GM]}}
                \texttt{PASIŪLYMAS: \{'A43', 'A67', 'B42', 'B19', 'C61', 'C67', 'C73', 'A18', 'C08', 'B36', 'C42'\}} \\
\\ 
\texttt{ARGUMENTAS: \{'Šis rinkinis apima visus svarbiausius man daiktus: A43, C42, C10 (keičiu C73, jei reikia, bet čia palieku C73 kaip papildomą vertę, jei tilpsta), B42, A67. Patikrinkime pastangas: A43(552) + A67(219) + B42(275) + B19(341) + C61(347) + C67(349) + C73(519) + A18(377) + C08(171) + B36(262) + C42(117) = 3529. Oh, viršija. Pataisykime. Pašalinti C73 (519). Tada suma: 3529 {-} 519 = 3010. Tai talpina C10? Ne, C10 (862) nepajudins. Geriau įtraukti C10. Pakeisti C73 (519) į C10 (862) + 343 skirtumas. 3010 + 343 = 3353. Tai telpa į 3358. Taigi rinkinys: A43, A67, B42, B19, C61, C67, C10, A18, C08, B36, C42. Patikrinimas: 552+219+275+341+347+349+862+377+171+262+117 = 3872? Aš suskaičiavau neteisingai. Leiskite perskaičiuoti.} \\
\texttt{A43: 552} \\
\texttt{A67: 219} \\
\texttt{B42: 275} \\
\texttt{B19: 341} \\
\texttt{C61: 347} \\
\texttt{C67: 349} \\
\texttt{C10: 862} \\
\texttt{A18: 377} \\
\texttt{C08: 171} \\
\texttt{B36: 262} \\
\texttt{C42: 117} \\
\texttt{Suma: 552+219+275+341+347+349+862+377+171+262+117 = 3872. Tai daug daugiau nei 3358.} \\
\\ 
\texttt{Reikia sumažinti.} \\
\texttt{Didžiausios pastangos: C10 (862), A43 (552), A18 (377), B19 (341), C61 (347), C67 (349).} \\
\texttt{Bendras limitas 3358.} \\
\texttt{Jei palikime C10, A43, C42, B42, A67, C08.} \\
\texttt{862 + 552 + 117 + 275 + 219 + 171 = 2196.} \\
\texttt{Likutis: 3358 {-} 2196 = 1162.} \\
\texttt{Galime įdėti: B19 (341), C61 (347), C67 (349). 341+347+349 = 1037.} \\
\texttt{Likutis: 1162 {-} 1037 = 125.} \\
\texttt{Galime įdėti B36 (262)? Ne. A18 (377)? Ne.} \\
\texttt{Taigi variantas: C10, A43, C42, B42, A67, C08, B19, C61, C67.} \\
\texttt{Suma: 862+552+117+275+219+171+341+347+349 = 3233. Tai telpa.} \\
\texttt{Ar tai geriau?} \\
\texttt{Mano svarba: C10(780)+A43(868)+C42(808)+B42(783)+A67(822)+C08(138)+B19(262)+C61(668)+C67(389) = 5668.} \\
\\ 
\texttt{Jūsų siūlymas (po pataisymo su C10):} \\
\texttt{C10, A43, A67, B42, B19, C61, C67, A18, C08, B36, C42.} \\
\texttt{Pastangos: 862+552+219+275+341+347+349+377+171+262+117 = 3872. Viršija.} \\
\texttt{Turime išmesti.} \\
\texttt{Išmetame A18 (377). Likutis 3495.} \\
\texttt{Išmetame B36 (262). Likutis 3233.} \\
\texttt{Tai yra tas pats rinkinys, kurį aš tik ką apskaičiavau, bet su C10, o ne C73.} \\
\texttt{Taigi: C10, A43, A67, B42, B19, C61, C67, C08, B36? Ne, B36 netelpa jei jau išmetėme.} \\
\texttt{Palikime C10, A43, A67, B42, B19, C61, C67, C08, C42.} \\
\texttt{Suma: 3233.} \\
\texttt{Ar galiu įdėti dar ką nors? Likutis 3358 {-} 3233 = 125. Mažiausias daiktas C08(171) jau yra. Kiti didesni.} \\
\texttt{Tai yra optimalus mano rinkinys?} \\
\texttt{Patikrinkime ar galiu pakeisti mažesnius į didesnius svarbos/kainos santykiu?} \\
\texttt{C42 (117, svarba 808) {-} labai geras.} \\
\texttt{C08 (171, svarba 138) {-} blogas santykis.} \\
\texttt{Galima pakeisti C08 į kitą daiktą, jei telpa?} \\
\texttt{Laisva vieta 125. C08 sveria 171. Negalima keisti į didesnį.} \\
\texttt{Galima išmesti C08 ir likti su 125 laisvos vietos.} \\
\texttt{Ar galiu įdėti B36 (262)? Ne.} \\
\\ 
\texttt{Taigi siūlau rinkinį be C08, jei jis man mažai vertas?} \\
\texttt{Bet C08 yra 171 pastangos. Jei jį išmesiu, turiu 3233 {-} 171 = 3062 pastangų ir 296 laisvos vietos.} \\
\texttt{Galėčiau įdėti B36 (262)? Taip.} \\
\texttt{Tada rinkinys: C10, A43, A67, B42, B19, C61, C67, C42, B36.} \\
\texttt{Suma: 3062 + 262 = 3324.} \\
\texttt{Mano svarba: 5668 {-} 138 (C08) + 461 (B36) = 5991.} \\
\texttt{Tai geriau!} \\
\\ 
\texttt{Taigi finalinis pasiūlymas:} \\
\texttt{\{'C10', 'A43', 'A67', 'B42', 'B19', 'C61', 'C67', 'C42', 'B36'\}} \\
\texttt{Pastangos: 862+552+219+275+341+347+349+117+262 = 3324 <= 3358.} \\
\texttt{Svarba: 780+868+822+783+262+66} \\
            }
        }
    }
    & & \\ \\

    \theutterance \stepcounter{utterance}  
    & & \multicolumn{4}{p{0.6\linewidth}}{
        \cellcolor[rgb]{0.9,0.9,0.9}{
            \makecell[{{p{\linewidth}}}]{
                \texttt{\tiny{[GM$\rangle$P2]}}
                \texttt{PASIŪLYMAS: \{'A43', 'A67', 'B42', 'B19', 'C61', 'C67', 'C73', 'A18', 'C08', 'B36', 'C42'\}} \\
\\ 
\texttt{ARGUMENTAS: \{'Šis rinkinis apima visus svarbiausius man daiktus: A43, C42, C10 (keičiu C73, jei reikia, bet čia palieku C73 kaip papildomą vertę, jei tilpsta), B42, A67. Patikrinkime pastangas: A43(552) + A67(219) + B42(275) + B19(341) + C61(347) + C67(349) + C73(519) + A18(377) + C08(171) + B36(262) + C42(117) = 3529. Oh, viršija. Pataisykime. Pašalinti C73 (519). Tada suma: 3529 {-} 519 = 3010. Tai talpina C10? Ne, C10 (862) nepajudins. Geriau įtraukti C10. Pakeisti C73 (519) į C10 (862) + 343 skirtumas. 3010 + 343 = 3353. Tai telpa į 3358. Taigi rinkinys: A43, A67, B42, B19, C61, C67, C10, A18, C08, B36, C42. Patikrinimas: 552+219+275+341+347+349+862+377+171+262+117 = 3872? Aš suskaičiavau neteisingai. Leiskite perskaičiuoti.} \\
\texttt{A43: 552} \\
\texttt{A67: 219} \\
\texttt{B42: 275} \\
\texttt{B19: 341} \\
\texttt{C61: 347} \\
\texttt{C67: 349} \\
\texttt{C10: 862} \\
\texttt{A18: 377} \\
\texttt{C08: 171} \\
\texttt{B36: 262} \\
\texttt{C42: 117} \\
\texttt{Suma: 552+219+275+341+347+349+862+377+171+262+117 = 3872. Tai daug daugiau nei 3358.} \\
\\ 
\texttt{Reikia sumažinti.} \\
\texttt{Didžiausios pastangos: C10 (862), A43 (552), A18 (377), B19 (341), C61 (347), C67 (349).} \\
\texttt{Bendras limitas 3358.} \\
\texttt{Jei palikime C10, A43, C42, B42, A67, C08.} \\
\texttt{862 + 552 + 117 + 275 + 219 + 171 = 2196.} \\
\texttt{Likutis: 3358 {-} 2196 = 1162.} \\
\texttt{Galime įdėti: B19 (341), C61 (347), C67 (349). 341+347+349 = 1037.} \\
\texttt{Likutis: 1162 {-} 1037 = 125.} \\
\texttt{Galime įdėti B36 (262)? Ne. A18 (377)? Ne.} \\
\texttt{Taigi variantas: C10, A43, C42, B42, A67, C08, B19, C61, C67.} \\
\texttt{Suma: 862+552+117+275+219+171+341+347+349 = 3233. Tai telpa.} \\
\texttt{Ar tai geriau?} \\
\texttt{Mano svarba: C10(780)+A43(868)+C42(808)+B42(783)+A67(822)+C08(138)+B19(262)+C61(668)+C67(389) = 5668.} \\
\\ 
\texttt{Jūsų siūlymas (po pataisymo su C10):} \\
\texttt{C10, A43, A67, B42, B19, C61, C67, A18, C08, B36, C42.} \\
\texttt{Pastangos: 862+552+219+275+341+347+349+377+171+262+117 = 3872. Viršija.} \\
\texttt{Turime išmesti.} \\
\texttt{Išmetame A18 (377). Likutis 3495.} \\
\texttt{Išmetame B36 (262). Likutis 3233.} \\
\texttt{Tai yra tas pats rinkinys, kurį aš tik ką apskaičiavau, bet su C10, o ne C73.} \\
\texttt{Taigi: C10, A43, A67, B42, B19, C61, C67, C08, B36? Ne, B36 netelpa jei jau išmetėme.} \\
\texttt{Palikime C10, A43, A67, B42, B19, C61, C67, C08, C42.} \\
\texttt{Suma: 3233.} \\
\texttt{Ar galiu įdėti dar ką nors? Likutis 3358 {-} 3233 = 125. Mažiausias daiktas C08(171) jau yra. Kiti didesni.} \\
\texttt{Tai yra optimalus mano rinkinys?} \\
\texttt{Patikrinkime ar galiu pakeisti mažesnius į didesnius svarbos/kainos santykiu?} \\
\texttt{C42 (117, svarba 808) {-} labai geras.} \\
\texttt{C08 (171, svarba 138) {-} blogas santykis.} \\
\texttt{Galima pakeisti C08 į kitą daiktą, jei telpa?} \\
\texttt{Laisva vieta 125. C08 sveria 171. Negalima keisti į didesnį.} \\
\texttt{Galima išmesti C08 ir likti su 125 laisvos vietos.} \\
\texttt{Ar galiu įdėti B36 (262)? Ne.} \\
\\ 
\texttt{Taigi siūlau rinkinį be C08, jei jis man mažai vertas?} \\
\texttt{Bet C08 yra 171 pastangos. Jei jį išmesiu, turiu 3233 {-} 171 = 3062 pastangų ir 296 laisvos vietos.} \\
\texttt{Galėčiau įdėti B36 (262)? Taip.} \\
\texttt{Tada rinkinys: C10, A43, A67, B42, B19, C61, C67, C42, B36.} \\
\texttt{Suma: 3062 + 262 = 3324.} \\
\texttt{Mano svarba: 5668 {-} 138 (C08) + 461 (B36) = 5991.} \\
\texttt{Tai geriau!} \\
\\ 
\texttt{Taigi finalinis pasiūlymas:} \\
\texttt{\{'C10', 'A43', 'A67', 'B42', 'B19', 'C61', 'C67', 'C42', 'B36'\}} \\
\texttt{Pastangos: 862+552+219+275+341+347+349+117+262 = 3324 <= 3358.} \\
\texttt{Svarba: 780+868+822+783+262+66} \\
            }
        }
    }
    & \\ \\

    \theutterance \stepcounter{utterance}  
    & & & \multicolumn{4}{p{0.6\linewidth}}{
        \cellcolor[rgb]{0.9,0.9,0.9}{
            \makecell[{{p{\linewidth}}}]{
                \texttt{\tiny{[GM$\langle$P2]}}
                \texttt{SUTINKU: \{'C10', 'A43', 'A67', 'B42', 'B19', 'C61', 'C67', 'C42', 'B36'\}} \\
            }
        }
    }
     \\ \\

    \theutterance \stepcounter{utterance}  
    & & \multicolumn{4}{p{0.6\linewidth}}{
        \cellcolor[rgb]{0.9,0.9,0.9}{
            \makecell[{{p{\linewidth}}}]{
                \texttt{\tiny{[GM$\rangle$P2]}}
                \texttt{Sutikote su pasiūlymu, kuris nėra aktyvus. Pasiūlymai yra aktyvūs tik tuo atveju, jei kitas žaidėjas juos užregistravo naudodamas PASIŪLYMAS ir jūs jų nedeaktyvavote naudodami ATSISAKAU. Bandykite dar kartą.} \\
            }
        }
    }
    & \\ \\

    \theutterance \stepcounter{utterance}  
    & & & \multicolumn{4}{p{0.6\linewidth}}{
        \cellcolor[rgb]{0.9,0.9,0.9}{
            \makecell[{{p{\linewidth}}}]{
                \texttt{\tiny{[GM$\langle$P2]}}
                \texttt{PASIŪLYMAS: \{'A43', 'A67', 'B42', 'B19', 'C61', 'C67', 'C10', 'C08', 'B36', 'C42'\}} \\
            }
        }
    }
     \\ \\

    \theutterance \stepcounter{utterance}  
    & & \multicolumn{4}{p{0.6\linewidth}}{
        \cellcolor[rgb]{0.9,0.9,0.9}{
            \makecell[{{p{\linewidth}}}]{
                \texttt{\tiny{[P1$\langle$GM]}}
                \texttt{PASIŪLYMAS: \{'A43', 'A67', 'B42', 'B19', 'C61', 'C67', 'C10', 'C08', 'B36', 'C42'\}} \\
            }
        }
    }
    & \\ \\

    \theutterance \stepcounter{utterance}  
    & \multicolumn{4}{p{0.6\linewidth}}{
        \cellcolor[rgb]{0.9,0.9,0.9}{
            \makecell[{{p{\linewidth}}}]{
                \texttt{\tiny{[P1$\rangle$GM]}}
                \texttt{SUTINKU: \{'A43', 'A67', 'B42', 'B19', 'C61', 'C67', 'C10', 'C08', 'B36', 'C42'\}} \\
            }
        }
    }
    & & \\ \\

    \theutterance \stepcounter{utterance}  
    & & & \multicolumn{2}{p{0.3\linewidth}}{
        \cellcolor[rgb]{0.9,0.9,0.9}{
            \makecell[{{p{\linewidth}}}]{
                \texttt{\tiny{[GM$|$GM]}}
                \texttt{game successful} \\
            }
        }
    }
    & & \\ \\

    \theutterance \stepcounter{utterance}  
    & & & \multicolumn{2}{p{0.3\linewidth}}{
        \cellcolor[rgb]{0.9,0.9,0.9}{
            \makecell[{{p{\linewidth}}}]{
                \texttt{\tiny{[GM$|$GM]}}
                \texttt{end game} \\
            }
        }
    }
    & & \\ \\

\end{supertabular}
}

\end{document}
